\section{Introduction}
\label{sec:introduction}

\subsection{Transient noise and why its shape matters}

The Advanced LIGO detectors \citep{AdvancedLIGO2015} reach strain sensitivities
of order $10^{-23}\,\mathrm{Hz}^{-1/2}$, and at that level the data are not
Gaussian. Short bursts of excess power --- \emph{glitches} --- occur at rates
of order one per minute, and they matter for two reasons: they reduce the
livetime available for astrophysical search, and they can mimic or corrupt real
signals \citep{Davis2021}. Detector characterisation exists to find them,
explain them, and where possible remove their cause from the instrument.

The productive question about a glitch is rarely ``how loud was it'' but
``what shape was it''. Glitch classes are defined morphologically, in the
time--frequency plane, and the morphology is what points back to a mechanism.
The Gravity Spy project \citep{Zevin2017,Glanzer2023} formalised this by
combining citizen-science labelling with a convolutional classifier over
$22$ morphological classes, and its classification release for the first three
observing runs is now the standard labelled dataset for this kind of work.

\subsection{Scattered light}

Among those classes, \emph{scattered light} is unusual in that its physical
mechanism is understood in closed form.

A small fraction of the circulating beam leaves the main path, strikes a
nearby surface --- a baffle, a suspended optic, a viewport --- and rejoins the
beam. Having traversed the extra path twice, the recombining field has
acquired a phase
%
\begin{equation}
  \phi(t) = \frac{4\pi\,x(t)}{\lambda},
  \label{eq:phase}
\end{equation}
%
where $x$ is the displacement of the scattering surface along the beam and
$\lambda = \SI{1064}{\nano\metre}$ is the Advanced LIGO laser wavelength. When
$\phi$ stays well below a radian the coupling is linear and the scattered light
simply adds a small amount of noise at the frequency of the motion. When $\phi$
exceeds a radian --- which requires only a fraction of a micron of
displacement --- the modulation is strongly non-linear, and the noise appears
at the instantaneous frequency
%
\begin{equation}
  f_1(t) = \frac{1}{2\pi}\left|\frac{\mathrm{d}\phi}{\mathrm{d}t}\right|
         = \frac{2}{\lambda}\,|v(t)|,
  \label{eq:f1}
\end{equation}
%
with $v = \mathrm{d}x/\mathrm{d}t$. Numerically,
$2/\lambda = \SI{1.88e6}{\hertz\per(\metre\per\second)}$: every micron per
second of surface velocity produces \SI{1.88}{\hertz} of noise frequency.

Light that makes $n$ round trips between the two surfaces before rejoining
accumulates $n$ times the phase, so the noise appears simultaneously at every
harmonic
%
\begin{equation}
  f_n(t) = n\,\frac{2}{\lambda}\,|v(t)|, \qquad n = 1, 2, 3, \ldots
  \label{eq:fn}
\end{equation}
%
with amplitudes falling off as the bounce order rises. In a Q-scan this
produces the characteristic picture: a comb of smooth, harmonically spaced
ridges, all peaking at the same instant.

Two consequences of \cref{eq:fn} carry the whole analysis.

\paragraph{The arch period measures the driving motion.} Because $f_n$ depends
on $|v|$ rather than on $v$, it reaches a maximum twice per cycle of the
motion. The observed arch repetition period is therefore \emph{half} the
period of the disturbance,
%
\begin{equation}
  \fdrive = \frac{1}{2\,\Tarch}.
  \label{eq:fdrive}
\end{equation}
%
A measured $\Tarch$ of \SI{3.3}{\second} implies a driver at
\SI{0.15}{\hertz}, squarely in the microseismic band; a measured
\SI{0.25}{\second} implies \SI{2}{\hertz}, in the anthropogenic band. The
morphology alone, with no auxiliary data at all, identifies the class of
disturbance.

\paragraph{The harmonic spacing measures the fundamental.} At the apex of an
arch train, all harmonics peak together, so a frequency slice through the apex
is a comb with spacing $f_1^{\max} = 2 v_{\mathrm{peak}}/\lambda$. The spacing
is thus a direct readout of the peak surface velocity, again from strain alone.

\subsection{What is already known}

\citet{Soni2021Scattering} characterised scattering throughout O3 and
identified two distinct populations at LIGO: a \emph{slow} component driven by
the \SIrange{0.1}{0.3}{\hertz} microseism, whose arches recur every few
seconds and reach a few tens of hertz, and a \emph{fast} component driven by
\SIrange{1}{6}{\hertz} motion, whose arches recur several times per second.
They traced the dominant O3 slow-scattering coupling to relative motion of a
few microns between the main suspension chain and its reaction chain, and
demonstrated mitigation by damping that relative motion. A witness-channel
algorithm predicting when scattering should be present is a central result of
that work, and \cref{eq:fn} is its basis.

Related work has established the surrounding picture: environmental coupling
paths in general \citep{Nguyen2021}, scattering from end benches at Virgo
\citep{Was2021,Accadia2010}, and the O3 data-quality products
\citep{Davis2021}. For the fourth run, \citet{Capote2025} document the
instrument performance and \citet{Soni2024O4} the detector characterisation of
its first half. Most directly relevant, \citet{Nandi2026} report on scattered
light at Livingston during O4 and again find two populations, one traced to
microseismic ground motion and mitigated by baffle installation, the other to
vertical ground motion coupling through a vacuum chamber.

That body of work is, however, built on internal LIGO data: full auxiliary
channel sets, engineering records, and access to the instrument. The present
study asks a narrower and quite different question.

\subsection{This work}

\emph{How much of the scattered-light story can be reconstructed from public
data alone, and does the answer change between O3 and O4a?}

The question is worth asking for two reasons. It is a genuine test of what the
open data releases support, which is useful to anyone outside the
collaboration. And it forces every inference to be made from quantities that
can be independently checked, which is a discipline that internal-data
analyses do not always impose.

The analysis proceeds in five steps.

\begin{enumerate}
  \item A labelled O3 sample of scattered-light triggers is drawn from the
        Gravity Spy release \citep{GravitySpyData2021}, together with a
        control sample of other glitch classes.
  \item Each trigger is Q-transformed \citep{Chatterji2004} and reduced to a
        twenty-element morphology feature vector, described in
        \cref{sec:methods}, whose components are the physical quantities of
        \crefrange{eq:fn}{eq:fdrive} rather than generic image statistics.
  \item The O3 sample is tied to physical cause using the public auxiliary
        channels: the hourly trigger rate is rank-correlated against
        seismometer band-limited RMS, and the predictor of \cref{eq:fn} is
        fitted trigger by trigger to a suspension witness.
  \item A classifier trained on the O3 features is applied to a
        livetime-limited blind search of O4a strain.
  \item The two populations are compared in rate, in slow/fast composition,
        and feature by feature.
\end{enumerate}

Steps 3 and 4 are shaped by two properties of the public releases that were
established before any code was written and that are reported here as results
in their own right, because they are not stated plainly anywhere else:

\begin{itemize}
  \item \textbf{No Gravity Spy classification catalogue exists for O4a.} The
        Zenodo release \citep{GravitySpyData2021} covers O1 through O3b. The
        O4a Gravity Spy model is described by \citet{Wu2025} but no catalogue
        accompanies it. Independent confirmation comes from
        \citet{Cirfeta2026}, who work on O4a without labels and use the O3b
        Gravity Spy set as their reference index --- exactly the workaround
        that O4a labels would render unnecessary.
  \item \textbf{The O4a auxiliary release contains no motion witness.} Its
        $25$ channels \citep{GWOSCO4Aux} are calibration lines, input
        mode-cleaner wavefront sensors, photon-calibrator readouts, mains
        monitors and parametric-instability monitors. There is no
        seismometer, accelerometer, microphone or suspension channel, so
        neither the rate correlation nor the per-trigger predictor can be
        computed for O4a from public data.
\end{itemize}

The O3 auxiliary release, by contrast, proved richer than expected:
$86$ full-rate channels and $402$ second-trend channels
\citep{GWOSCO3Aux}, including the complete ground-seismometer band-limited RMS
family across the earthquake, microseismic and anthropogenic bands at both
sites, and --- crucially --- \code{L1:SUS-OM3\_M1\_DAMP\_L\_IN1\_DQ} at
\SI{256}{\hertz}, a genuine optic-motion witness and the only full-rate one in
either release.

\Cref{sec:data} sets out the data and verifies its availability;
\cref{sec:methods} the analysis; \cref{sec:validation} the validation against
simulated injections with known parameters; \cref{sec:results} the results;
\cref{sec:discussion} their interpretation and limits.
